\documentclass[oneside, a4paper, onecolumn, 11pt]{article}

% Change this: Customize the title, author, advisor, abstract
\newcommand{\thesistitle}[0]{Optimizing the implementation of Witt vectors in the SageMath software}
\newcommand{\authorname}[0]{Yonas Grossard-Amin}

\newcommand{\supervisor}[0]{Rub\'en Mu\~noz-\relax-Bertrand}
\newcommand{\supervisorinstitution}[0]{Inria Saclay}

\newcommand{\abstracttext}[0]{%
Later
}

\usepackage[
  left=2cm,top=2.0cm,bottom=2.0cm,right=2cm,
  headheight=17pt, % as per the warning by fancyhdr
  includehead,includefoot,
  heightrounded, % to avoid spurious underfull messages
]{geometry}


\usepackage[T1]{fontenc}
\usepackage{amstext}
\usepackage{amsmath}
\usepackage{amssymb}
\usepackage{url}
\usepackage{graphicx}
\usepackage{wrapfig}
\usepackage{enumerate}
\usepackage{paralist}
\usepackage{xspace}
\usepackage{color}
\usepackage{times}
\usepackage[colorlinks,linkcolor=blue]{hyperref}
\usepackage[colorinlistoftodos,prependcaption,textsize=normal]{todonotes}
\usepackage{pdfpages}
\usepackage{fancyhdr} %% For changing headers and footers

\usepackage{titling}
\usepackage[nottoc,numbib]{tocbibind}

\usepackage{amsthm}
\newtheorem{theorem}{Theorem}
\newtheorem{lemma}{Lemma}
\newtheorem{proposition}{Proposition}
\newtheorem{corollary}{Corollary}
\theoremstyle{definition}
\newtheorem{definition}{Definition}[section]
\usepackage{algorithm}
\usepackage{algpseudocode}
\usepackage{float}

%% \predate{}
%% \postdate{}
%% \date{}
%% \author{\authorname}


\begin{document}

%\title{\thesistitle}

%\maketitle

% Max 10 lines.
%\noindent \paragraph*{Abstract}
%\abstract

\hspace{0pt}
\vfill

\begin{center}

\includegraphics[width=0.3\textwidth]{logo-EP-vertical}

\vspace*{2em}
%
{\large
\textbf{\'Ecole Polytechnique}

\vspace*{1em}
\textit{BACHELOR THESIS IN COMPUTER SCIENCE}


\vspace*{3em}
{\Huge \textbf{\thesistitle}}
\vspace*{3em}



\textit{Author:}

\vspace*{1em}
\authorname{}, \'Ecole Polytechnique

\vspace*{2em}
%
{\textit{Advisor:}}

\vspace*{1em}
\supervisor{}, \supervisorinstitution{}
}

\vspace*{2em}
\textit{Academic year 2025/2026}

\end{center}

\vfill
\hspace{0pt}

\newpage

\vfill
\noindent\textbf{Abstract}\\[1em]
%
\fbox{
\parbox{\textwidth}{
\abstracttext{}
}
}
\vfill


\newpage

% Setting up the header
\pagestyle{fancy}
%\renewcommand{\headrulewidth}{0pt} % Remove line at top
%\renewcommand{\headrulewidth}{0.4pt}% Default \headrulewidth is 0.4pt
\lhead{\authorname}
%\chead{\acronym}
\rhead{\thesistitle}



\newpage
\tableofcontents
\newpage

%\pagenumbering{arabic}

\section{Introduction}

\subsection{Motivation}
Witt vector rings are algebraic objects with multiple useful properties such as the possibility of \emph{lifting} rings from characteristic $p$, with $p$ a prime, to characteristic $0$, as well as being an alternative way of constructing \emph{p-adic numbers}. 

Computationally we use truncated Witt vectors, the finite counterpart to Witt vectors, as it has the same properties in the finite case. The problem is that the construction of the Witt vector polynomials are of exponential complexity, both in $p$ and in the length of the truncated Witt vectors, called \emph{precision}.

So caching etc is important.

\section{Witt Vector Rings}
\subsection{Mathematical background}\label{subsec:background}
\fbox{\parbox{\textwidth}{
These results are shown before the construction of Witt vectors, as they are important for said construction but can confuse the reader if intertwined with the explanation. These results can be considered as prerequisites for the understanding of the construction of Witt vectors.
}}
\newline

In the following, let $p$ be a fixed prime number. 

It is common knowledge that 
\begin{proposition}
  $\mathbb{F}_p = \mathbb{Z}/p\mathbb{Z} = \mathbb{Z}/(p)$ endowed with the usual addition and multiplication modulo $p$ is a field with a finite amount of elements.
\end{proposition} 
This can easily be checked, and can also be proven using the fact that the ideal $(p)$ is maximal, and that the quotient of a ring with a maximal ideal is a field.  \textbf{Maybe proof in appendix ? not sure its necessary}

Now recall the construction of the Frobenius map:
\begin{definition}
  Let $A$ be a ring of characteristic $p$. \[F: A \to A\] defined by \[F(x) = x^p\] is called the Frobenius map.
\end{definition} 
\begin{proposition}
  The Frobenius map is a ring homomorphism.
\end{proposition}
\begin{proof}
  For all $x,y\in A$, one has:
  \[F(xy) = (xy)^p = x^p y^p = F(x)F(y).\]
  And also \[F(x+y) = (x+y)^p.\]
  Where we can utilize the fact that in modulo $p$, for all $a,b\in \mathbb{F}_p$ we have $(a+b)^p=a^p+b^p$. In fact, \[(a+b)^p = \sum_{i=0}^{p}\binom{p}{i}a^ib^{p-i} = a^p+b^p+\sum_{i=1}^{p-1}\binom{p}{i}a^ib^{p-i}.\]
  And the binomial factor $\binom{p}{i}$ is a multiple of $p$ for all $i\in\{1,\ldots,p-1\}$, so it is zero in modulo $p$, so
  \[(x+y)^p = x^p + y^p.\]
  Therefore \[F(x+y)=F(x)+F(y).\]
\end{proof}
\begin{proposition}
  In the case $A=\mathbb{F}_p$ the Frobenius map is the identity.
\end{proposition}
\begin{proof}
  Let $a\in \mathbb{F}_p$. If $a=0$ then the identity holds. Else, $a\in\mathbb{F}_p^{\times}= \mathbb{F}_p\backslash\{0\}$
  and then, by Fermat's little theorem, \[a^p\equiv a \pmod p.\]
  Since $\mathbb{F}_p = \mathbb{Z}/p\mathbb{Z}$, congruence modulo $p$ becomes equality, and the Frobenius map is indeed the identity.
\end{proof}
From this, if $a\in\mathbb{F}_p^{\times}$ we get that $a^{p-1}\equiv1\pmod p$ by multiplying both sides by $a^{-1}$. This means that each non-zero element of $\mathbb{F}_p$ is a $(p-1)$th root of unity.

\begin{definition}
  Let $x\in\mathbb{Z}_p$ then \[x=\sum_{i=0}^{\infty}a_ip^i \text{  with  } a_i\in\{0,..,p-1\}.\]
\end{definition}
\begin{definition}
For a nonzero integer $x\in\mathbb Z$, the \emph{$p$-adic valuation} $v_p(x)$ is defined as the largest integer $k\ge0$ such that
\[p^k \mid x.\]
We set $v_p(0)=\infty$.
\end{definition}
\begin{definition}
The \emph{$p$-adic absolute value} on $\mathbb Z$ is defined by \[|x|_p = p^{-v_p(x)}, \qquad |0|_p=0.\]
\end{definition}

This absolute value satisfies\[|xy|_p = |x|_p |y|_p\]

and the inequality \[|x+y|_p \le \max(|x|_p,|y|_p).\]

\begin{definition}
A sequence $(x_n)$ of integers is called \emph{$p$-adically Cauchy} if for every $k\ge1$ there exists $N$ such that
\[x_m-x_n \in p^k\mathbb Z\qquad\text{for all } m,n\ge N.\]
\end{definition}
The ring of \emph{$p$-adic integers} $\mathbb Z_p$ is defined as the completion of $\mathbb Z$ with respect to this notion of Cauchy sequence. Elements of $\mathbb Z_p$ can therefore be represented as limits of $p$-adically Cauchy sequences of integers.
The ring $\mathbb Z_p$ is complete with respect to the $p$-adic topology.

A sequence $(x_n)$ in $\mathbb Z_p$ converges to $0$ if and only if
\[v_p(x_n) \to \infty.\]

In particular,\[p^n \to 0\qquad\text{in } \mathbb Z_p,\] since $v_p(p^n)=n$.

More generally, if $x_n\in\mathbb Z_p$, then \[p^n x_n \to 0.\]

Finally, polynomial maps are continuous for the $p$-adic topology. In particular the map \[x \mapsto x^p\] is continuous on $\mathbb Z_p$.

\subsection{Construction from p-adic expansions}

There are multiple ways of introducing Witt vectors. Instead of just stating the definitions, we will try to follow the thought process which leads to the creation of these definitions. 

We will first show how the idea of Witt vectors come to light naturally, then how one comes to the idea of ghost components, finally we will show how the operations result in sum and product polynomials.
\newline
\newline
In the following, let $p$ be a fixed prime number. 

Every $p$-adic integer $s \in \mathbb{Z}_p$ admits a unique expansion of the form \[s=\sum_{i=0}^{\infty}a_ip^i \text{  with  } a_i\in\{0,..,p-1\}.\]

Since the set $\{0,1,\ldots,p-1\}$ is bijective to $\mathbb{F}_p$, each $p$-adic number can be uniquely represented by a sequence $(a_0,a_1,a_2,\ldots)$ with $a_i \in \mathbb{F}_p$.

A natural first idea is to try to reconstruct a $p$-adic integer directly from a sequence
\[(a_0,a_1,a_2,\ldots)\in \mathbb{F}_p^{\mathbb N}\]
by forming the formal expression
\[a_0+a_1p+a_2p^2+\cdots.\]
However, this expression is not yet meaningful in $\mathbb{Z}_p$, because the coefficients $a_i$ lie in $\mathbb{F}_p$, whereas a $p$-adic expansion requires coefficients in $\mathbb{Z}_p$.

We therefore need a way to choose, for each element of $\mathbb{F}_p$, a distinguished representative in $\mathbb{Z}_p$. In other words, we seek a map
\[\chi:\mathbb{F}_p\to\mathbb{Z}_p\]
which is a section of the reduction map
\[\pi:\mathbb{Z}_p\to\mathbb{F}_p,\qquad \pi(x)=x \bmod p.\]
By definition, this means that
\[\pi(\chi(a))=a\qquad \text{for all } a\in\mathbb{F}_p.\]
Equivalently, we require
\[\chi(a)\equiv a \pmod p.\]

Then \[\pi(\chi(a))=a \qquad\text{for all } a\in\mathbb{F}_p.\]
Equivalently, we require \[ \chi(a)\equiv a \pmod p.\]

In order for this lift to interact well with the arithmetic in characteristic $p$, it is natural to impose the additional condition that $\chi(a)$ be fixed by Frobenius: \[\chi(a)^p=\chi(a).\]
Thus we are now looking for elements $x\in\mathbb{Z}_p$ satisfying $x^p=x$.

In fact such a lift exists, and it is unique. 

\begin{proposition}
For every $a\in\mathbb{F}_p$, there exists a unique element $\chi(a)\in\mathbb{Z}_p$ such that
\[\chi(a)\equiv a \pmod p \qquad \text{and} \qquad \chi(a)^p=\chi(a).\]
This map $\chi:\mathbb{F}_p\to\mathbb{Z}_p$ is called the Teichmüller lift.
\end{proposition}

\begin{proof}
First consider the case $a=0$. Then $x=0$ clearly satisfies
\[ x\equiv 0 \pmod p\qquad\text{and}\qquad x^p=x\]
and it is the unique such element. It is unique, since if $x\ne 0$ then $x=pk$ for some $k\in\mathbb{Z}_p$, and then $x^p=p^pk^p=pkp^{p-1}k^{p-1}=xp^{p-1}k^{p-1}\ne x$.


Now let $a\in\mathbb{F}_p^\times$. Choose its standard representative
$a\in\{1,\dots,p-1\}\subset\mathbb{Z}\subset\mathbb{Z}_p$.
Define a sequence $(x_n)_{n\ge0}$ in $\mathbb{Z}_p$ by
\[x_n:=a^{p^n}.\]

We claim that $(x_n)$ is a Cauchy sequence in $\mathbb{Z}_p$. Since
\[x_{n+1}=x_n^p,\]
and since $x_n\equiv a \pmod p$, we have
\[x_{n+1}=x_n^p\equiv x_n \pmod p.\]
More generally, applying Lemma \ref{equiv+i} repeatedly gives
\[x_{n+1}=x_n^p \equiv x_n \pmod{p^{n+1}}.\]
Thus
\[x_{n+1}-x_n\in p^{n+1}\mathbb{Z}_p,\]
so $(x_n)$ is Cauchy. Since $\mathbb{Z}_p$ is complete, there exists
$x\in\mathbb{Z}_p$ such that
\[x_n\to x.\]

Because $x_n\equiv a \pmod p$ for all $n$, passing to the limit gives
\[x\equiv a \pmod p.\]
Also, the map $t\mapsto t^p$ is continuous on $\mathbb{Z}_p$, so
\[x^p=\lim_{n\to\infty} x_n^p=\lim_{n\to\infty} x_{n+1}=x.\]
We therefore define $\chi(a):=x$.

This proves existence.

To prove uniqueness, let $x,y\in\mathbb{Z}_p$ satisfy
\[x\equiv a \pmod p,\qquad y\equiv a \pmod p,\qquad x^p=x,\qquad y^p=y.\]
Then $x\equiv y\pmod p$. By Lemma \ref{equiv+i},
\[x^{p^n}\equiv y^{p^n}\pmod{p^{n+1}}\qquad\text{for all } n\ge0.\]
Since $x^p=x$ and $y^p=y$, this becomes
\[x\equiv y\pmod{p^{n+1}}\qquad\text{for all } n\ge0.\]
Hence $x-y\in \bigcap_{n\ge1} p^n\mathbb{Z}_p=\{0\}$, so $x=y$.

Thus for every $a\in\mathbb{F}_p$ there exists a unique element
$\chi(a)\in\mathbb{Z}_p$ satisfying
\[\chi(a)\equiv a \pmod p\qquad\text{and}\qquad\chi(a)^p=\chi(a).\]
\end{proof}

Having chosen a canonical lift of each element of $\mathbb{F}_p$, we may now define a map
\[\phi:\mathbb{F}_p^{\mathbb N}\to \mathbb{Z}_p\]
by
\[\phi((a_0,a_1,a_2,\ldots))=\chi(a_0)+p\chi(a_1)+p^2\chi(a_2)+\cdots.\]
This is a well-defined element of $\mathbb{Z}_p$, as it is a sum of elements of $\mathbb{Z}_p$, and it expresses a $p$-adic integer in terms of Teichmüller representatives.

Moreover, every element of $\mathbb{Z}_p$ admits a unique expansion of this form, by construction of $\mathbb{Z}_p$. Therefore, $\phi$ is a bijection of sets
\[\phi:\mathbb{F}_p^{\mathbb N}\xrightarrow{\sim}\mathbb{Z}_p.\]

We now investigate whether such a reconstruction can be made into a ring-theoretic one.

Since $\mathbb{Z}_p$ is a ring, we need to endow $\mathbb{F}_p^{\mathbb{N}}$ with a ring structure that supports the arithmetic of $p$-adic integers. However this cannot be done using componentwise operations because of the carry arithmetic involved in addition and multiplication of $p$-adic integers. For example:
\[(1,0,0,\ldots)+(p-1,0,0,\ldots)=(0,0,0,\ldots)\] under componentwise addition in $\mathbb{F}_p^{\mathbb{N}}$ whereas \[1+(p-1)=p\] in $\mathbb{Z}_p$ corresponds to \[(0,1,0,\ldots)\ne (0,0,0,\ldots).\]

Thus the correct ring structure on $\mathbb{F}_p^{\mathbb{N}}$ cannot be the naive componentwise one. Instead, we will try to transpose the ring structure of $\mathbb{Z}_p$ to $\mathbb{F}_p^{\mathbb{N}}$ using $\phi$. But first, let us prove that $\phi$ is a bijection.

\begin{proposition}
The map
\[\phi:\mathbb F_p^{\mathbb N}\to \mathbb Z_p,\qquad (a_0,a_1,a_2,\ldots)\mapsto \sum_{i=0}^\infty p^i\chi(a_i)\]
is a bijection.
\end{proposition}

\begin{proof}
We first prove injectivity.

Let
\[a=(a_0,a_1,a_2,\ldots), \qquad b=(b_0,b_1,b_2,\ldots)\in\mathbb F_p^{\mathbb N}\]
and suppose that $\phi(a)=\phi(b)$. Then
\[\chi(a_0)+p\chi(a_1)+p^2\chi(a_2)+\cdots=\chi(b_0)+p\chi(b_1)+p^2\chi(b_2)+\cdots.\]
Reducing modulo $p$, we obtain
\[\chi(a_0)\equiv \chi(b_0)\pmod p.\]
Applying the reduction map and using that $\pi\circ\chi=\mathrm{id}_{\mathbb F_p}$, it follows that
\[a_0=b_0.\]

Subtracting the common first term and dividing by $p$, we get
\[\chi(a_1)+p\chi(a_2)+p^2\chi(a_3)+\cdots=\chi(b_1)+p\chi(b_2)+p^2\chi(b_3)+\cdots.\]
Reducing again modulo $p$, we obtain $a_1=b_1$.

Repeating this argument inductively shows that $a_i=b_i$ for all $i\ge 0$. Hence $a=b$, and $\phi$ is injective.

We now prove surjectivity.

Let $x\in\mathbb Z_p$. We will construct a sequence
\[(a_0,a_1,a_2,\ldots)\in\mathbb F_p^{\mathbb N}\]
such that
\[x=\sum_{i=0}^\infty p^i\chi(a_i).\]

Let $a_0:=\pi(x)\in\mathbb F_p$. Since $\pi(\chi(a_0))=a_0=\pi(x)$, we have
\[x-\chi(a_0)\in p\mathbb Z_p.\]
Hence there exists $x_1\in\mathbb Z_p$ such that
\[x=\chi(a_0)+px_1.\]

Next, let $a_1:=\pi(x_1)\in\mathbb F_p$. Again, since
\[\pi(\chi(a_1))=a_1=\pi(x_1),\]
we have
\[x_1-\chi(a_1)\in p\mathbb Z_p.\]
Thus there exists $x_2\in\mathbb Z_p$ such that
\[x_1=\chi(a_1)+px_2.\]
Substituting into the previous equality gives
\[x=\chi(a_0)+p\chi(a_1)+p^2x_2.\]

Continuing inductively, we construct sequences $(a_n)_{n\ge0}$ in $\mathbb F_p$
and $(x_n)_{n\ge0}$ in $\mathbb Z_p$ such that
\[x=\chi(a_0)+p\chi(a_1)+\cdots+p^n\chi(a_n)+p^{n+1}x_{n+1}\qquad\text{for all } n\ge0.\]

Now the remainder term $p^{n+1}x_{n+1}$ tends to $0$ in $\mathbb Z_p$ as
$n\to\infty$, because $x_{n+1}\in\mathbb Z_p$. Therefore
\[x=\sum_{i=0}^\infty p^i\chi(a_i).\]
Thus
\[x=\phi((a_0,a_1,a_2,\ldots)),\]
and so $\phi$ is surjective.

Therefore $\phi$ is bijective.
\end{proof}

\subsection{Ghost components}

Because the map $\phi$ provides a bijection between $\mathbb{F}_p^{\mathbb{N}}$ and $\mathbb{Z}_p$, we will try to endow $\mathbb{F}_p$ with a ring structure, using the already established one in $\mathbb{Z}_p$.
Let us try to define, for two sequences \[a=(a_0,a_1,a_2,\ldots), \qquad b=(b_0,b_1,b_2,\ldots), \qquad a,b\in\mathbb{F}_p^{\mathbb{N}}\]
the operations \[a \oplus b := \phi^{-1}(\phi(a)+\phi(b)), \qquad a \otimes b := \phi^{-1}(\phi(a)\phi(b)).\]

This indeed defines a ring structure on $\mathbb{F}_p^{\mathbb{N}}$ that is isomorphic to $\mathbb{Z}_p$. However, the resulting formulas for the coordinates of $a\oplus b$ and $a\otimes b$ are complicated due to the carry terms appearing in the arithmetic of $p$-adic expansions.

To understand this phenomenon more precisely, let \[x=\phi(a), \qquad y=\phi(b).\]
Then \[x=\chi(a_0)+p\chi(a_1)+p^2\chi(a_2)+\cdots\]
and \[y=\chi(b_0)+p\chi(b_1)+p^2\chi(b_2)+\cdots.\]

Now let $c\in\mathbb{F}_p^{\mathbb{N}}$ such that \[ x+y=\phi(c)=\chi(c_0)+p\chi(c_1)+p^2\chi(c_2)+\cdots.\]
This then gives the following congruences:
\begin{equation}\label{original_congruence}
\begin{aligned}
\chi(c_0) &\equiv \chi(a_0) + \chi(b_0) \pmod{p} \\
\chi(c_0) + \chi(c_1)p &\equiv \chi(a_0) + \chi(a_1)p + \chi(b_0) + \chi(b_1)p \pmod{p^2} \\
\chi(c_0) + \chi(c_1)p + \chi(c_2)p^2 &\equiv \chi(a_0) + \chi(a_1)p + \chi(a_2)p^2  + \chi(b_0) + \chi(b_1)p + \chi(b_2)p^2 \pmod{p^3} \\
&\vdots
\end{aligned}
\end{equation}
The first congruence directly implies that $c_0=a_0+b_0$.

Now we will prove the following lemma, which will help solve for $c_i$ for $i>1$.

\begin{lemma}\label{equiv+i}
For $x,y\in \mathbb{F}_p$, If $x\equiv y \pmod p$ then $x^{p^i}\equiv y^{p^i} \pmod {p^{i+1}}$ for every $i\in\mathbb{N}$.
\end{lemma}
\begin{proof}
  We will proceed by recursion on $i\in\mathbb{N}$, with the case $i=0$ being the hypothesis.
  Suppose now that for $i\in\mathbb{N}$ we have $x^{p^i}\equiv y^{p^i} \pmod {p^{i+1}}$. We can write, for some $\alpha \in\mathbb{F}_p $ \[x^{p^i}=y^{p^i}+\alpha p^{i+1}\implies x^{p^{i+1}}=\bigl(y^{p^i}+\alpha p^{i+1}\bigr)^p.\]
  Now expanding with Newton's binomial theorem, \[x^{p^{i+1}}=\bigl(y^{p^(i+1)}+\alpha p^{i+1}\bigr)^p=y^{p^i}+\alpha^p p^{p(i+1)}+\sum_{k=0}^{p-1}\binom{p}{k}y^{p^{i}(p-k)}\alpha^k p^{k(i+1)}.\]
  For $k\in\{1,\ldots,p-1\}$ we have \[\binom{p}{k}\mid p \implies \binom{p}{k}y^{p^{i}(p-k)}\alpha^k p^{k(i+1)}\mid p^{i+2}.\]
  And \[p^{p(i+1)}\mid p^{i+2} \text{ since }p\geq 2.\]
  We indeed get that $x^{p^{i+1}}\equiv y^{p^{i+1}} \pmod {p^{i+2}}$.
\end{proof}
Trying to find $c_1$ using the congruence above, we get \[\chi(c_1)p \equiv \chi(a_0) + \chi(a_1)p + \chi(b_0) + \chi(b_1)p - \chi(c_0) \pmod{p^2}.\]

Using that $\chi(c_0) \equiv \chi(a_0) + \chi(b_0) \pmod{p}$ we get that $\chi(a_0) + \chi(b_0) - \chi(c_0)  \mid p$ so we can divide by $p$ \begin{equation}\chi(c_1) \equiv  \chi(a_1) +  \chi(b_1) + \frac{\chi(a_0) + \chi(b_0) - \chi(c_0)}{p}\pmod{p}.\end{equation}\label{congruence}
Now using the lemma \ref{equiv+i} and the properties of the Teichmüller lift, we get that: \[\chi(c_0) \equiv \chi(a_0) + \chi(b_0) \pmod{p} \implies \chi(c_0)^p \equiv \bigl(\chi(a_0) + \chi(b_0)\bigr)^p \pmod{p^2}\]\[ \implies \chi(c_0) \equiv \bigl(\chi(a_0) + \chi(b_0)\bigr)^p \pmod{p^2} \]
We will expand Newton's binomial theorem $\bigl(\chi(a_0) + \chi(b_0)\bigr)^p=\chi(a_0)^p+\chi(b_0)^p+\sum_{k=1}^{p-1}\binom{p}{k}\chi(a_0)^k\chi(b_0)^{p-k}$. 
Now note that \[\binom{p}{k}=\frac{p!}{k!(p-k)!}\implies \binom{p}{k}\mid p \text{ and } \binom{p}{k}\nmid p^2 \text{ for } k\in\{1,\ldots,p-1\}.\]
So \[\chi(c_0) \equiv \chi(a_0)+\chi(b_0)+p(\chi(a_0)^{p-1}\chi(b_0)+\ldots+\chi(a_0)\chi(b_0)^{p-1}) \pmod{p^2}.\]
Reordering gets us \[\chi(a_0)+\chi(b_0)-\chi(c_0)\equiv -p(\chi(a_0)^{p-1}\chi(b_0)+\ldots+\chi(a_0)\chi(b_0)^{p-1})\pmod{p^2}.\]
Again, since  $\chi(a_0) + \chi(b_0) - \chi(c_0)  \mid p$, \[\frac{\chi(a_0) + \chi(b_0) - \chi(c_0)}{p}\equiv -(\chi(a_0)^{p-1}\chi(b_0)+\ldots+\chi(a_0)\chi(b_0)^{p-1})\pmod{p}\]
Now injecting this into \ref{congruence} gets us \[\chi(c_1) \equiv  \chi(a_1) +  \chi(b_1) -(\chi(a_0)^{p-1}\chi(b_0)+\ldots+\chi(a_0)\chi(b_0)^{p-1}) \pmod{p}\]
\[c_1 \equiv a_1 + b_1 -(a_0^{p-1}b_0+\ldots+a_0b_0^{p-1}) \pmod{p}.\]
So \[c_1 = a_1 + b_1 -(a_0^{p-1}b_0+\ldots+a_0b_0^{p-1})\] since we have elements in $\mathbb{F}_p$.

As you can see, this computation is quite heavy because of the \emph{carry term}, and it becomes even more complicated when computing more. One idea to make this computation easier, is to change the coordinate system into one that has easy addition. Looking back at \ref{original_congruence} and using the lemma \ref{equiv+i} we get:
\begin{align*}
\chi(c_0) &\equiv \chi(a_0) + \chi(b_0) \pmod{p} \\
\chi(c_0)^p + \chi(c_1)p &\equiv \chi(a_0)^p + \chi(a_1)p + \chi(b_0)^p + \chi(b_1)p \pmod{p^2} \\
\chi(c_0)^{p^2} + \chi(c_1)^pp + \chi(c_2)p^2 &\equiv \chi(a_0)^{p^2} + \chi(a_1)^pp + \chi(a_2)p^2  + \chi(b_0)^{p^2} + \chi(b_1)^pp + \chi(b_2)p^2 \pmod{p^3} \\
&\vdots
\end{align*}
From this, one could define:
\begin{definition}
For a sequence $a=(a_0,a_1,a_2,\ldots)\in\mathbb{F}_p^{\mathbb{N}}$, define its $n$-th ghost component by
\[w_n(a)=a_0^{p^n}+pa_1^{p^{n-1}}+p^2a_2^{p^{n-2}}+\cdots+p^n a_n.\]
The family $(w_0(a),w_1(a),w_2(a),\ldots)$ is called the sequence of ghost components of $a$.


The map
\[w:\mathbb{F}_p^{\mathbb N}\to \mathbb{Z}_p^{\mathbb N},\qquad a\mapsto (w_0(a),w_1(a),w_2(a),\ldots)\]
is called the ghost map.
\end{definition}
The name \emph{ghost components} reflects the fact that these are not the original coordinates of the Witt vector, but rather polynomial expressions in those coordinates. They are auxiliary quantities, invisible from the point of view of the original sequence, yet they encode the arithmetic in a much simpler way, allowing for easier computations.
With this notation, the congruences obtained above can be rewritten in the form
\[w_0(c)\equiv w_0(a)+w_0(b)\pmod p,\]
\[w_1(c)\equiv w_1(a)+w_1(b)\pmod{p^2},\]
\[w_2(c)\equiv w_2(a)+w_2(b)\pmod{p^3},\]
and so on. The ghost components behave additively modulo the appropriate power of $p$.
\subsection{Universal Witt polynomials}
Now that we have logically introduced the idea of \emph{ghost components} when trying to lift $\mathbb{F}_p$ into a characteristic $0$ ring, namely $\mathbb{Z}_p$, we will instead broaden this logic to any commutative ring $A$. To do so, we will first restate \ref{equiv+i} in $A$, then state and prove an auxiliary lemma to finally prove the uniqueness and existence of the sum and product polynomials, which define a ring structure on $W(A)=A^{\mathbb{N}}$.

\begin{lemma}\label{equiv_A}
  Let $A$ be a commutative ring. For $x,y\in A$, If $x\equiv y \pmod pA$ then $x^{p^i}\equiv y^{p^i} \pmod {p^{i+1}A}$ for every $i\in\mathbb{N}$.
\end{lemma}
\begin{proof}
  The proof is very similar to its counterpart for $A=\mathbb{F}_p$. It is in the appendix \ref{sec:appendix}.
\end{proof}

\begin{lemma}\label{lemma:frobenius_congruence}
Let $A$ be a commutative ring, and let $\varphi:A\to A$ be a ring endomorphism such that
\[\varphi(a)\equiv a^p \pmod{pA}\qquad \text{for all } a\in A.\]
Let $(u_n)_{n\geq 0}$ be a sequence of elements of $A$ satisfying
\[\varphi(u_n)\equiv u_{n+1}\pmod{p^{n+1}A}\qquad \text{for all } n\geq 0.\]
Then there exists a unique sequence $(z_n)_{n\geq 0}$ in $A$ such that
\[u_n=z_0^{p^n}+pz_1^{p^{n-1}}+\cdots+p^n z_n\qquad \text{for all } n\geq 0.\]
\end{lemma}

\begin{proof}

We prove existence and uniqueness simultaneously by induction on $n$.

For $n=0$, the required identity is simply
\[u_0=z_0,\]
so we must take
\[z_0=u_0.\]

Now suppose that $z_0,\dots,z_{n-1}$ have already been constructed so that
\[u_m=z_0^{p^m}+pz_1^{p^{m-1}}+\cdots+p^m z_m\qquad \text{for all } m\le n-1.\]
We now show that $z_n$ exists and is uniquely determined.

Consider the expression
\[N_n:=u_n-\bigl(z_0^{p^n}+pz_1^{p^{n-1}}+\cdots+p^{n-1}z_{n-1}^p\bigr).\]
We claim that $N_n$ is divisible by $p^n$ in $A$.

Indeed, from the induction hypothesis applied to $u_{n-1}$, we have
\[u_{n-1}=z_0^{p^{n-1}}+pz_1^{p^{n-2}}+\cdots+p^{n-1}z_{n-1}.\]
Applying $\varphi$, we get
\[\varphi(u_{n-1})=\varphi(z_0)^{p^{n-1}}+p\,\varphi(z_1)^{p^{n-2}}+\cdots+p^{n-1}\varphi(z_{n-1}).\]
Since $\varphi(a)\equiv a^p\pmod{pA}$ for every $a\in A$, one has
\[\varphi(z_i)=z_i^p+p\,b_i\]
for some $b_i\in A$.
Hence
\[\varphi(z_i)^{p^{n-1-i}}\equiv z_i^{p^{n-i}}\pmod{p^{\,n-i}A}.\]
Multiplying by $p^i$, we obtain
\[p^i\varphi(z_i)^{p^{n-1-i}} \equiv p^i z_i^{p^{n-i}}\pmod{p^n A}.\]
Summing over $i=0,\dots,n-1$, it follows that
\[\varphi(u_{n-1})\equiv z_0^{p^n}+pz_1^{p^{n-1}}+\cdots+p^{n-1}z_{n-1}^p\pmod{p^nA}.\]
But by hypothesis,
\[u_n\equiv \varphi(u_{n-1})\pmod{p^nA}.\]
Combining the two congruences yields
\[u_n\equiv z_0^{p^n}+pz_1^{p^{n-1}}+\cdots+p^{n-1}z_{n-1}^p\pmod{p^nA}.\]
Thus $N_n\in p^nA$, as claimed.

Therefore there exists $z_n\in A$ such that
\[N_n=p^n z_n,\]
that is,
\[u_n=z_0^{p^n}+pz_1^{p^{n-1}}+\cdots+p^{n-1}z_{n-1}^p+p^n z_n.\]
This proves existence.

For uniqueness, observe that once $z_0,\dots,z_{n-1}$ are fixed, the above identity determines $z_n$ uniquely, since
\[p^n z_n=u_n-\bigl(z_0^{p^n}+pz_1^{p^{n-1}}+\cdots+p^{n-1}z_{n-1}^p\bigr).\]
Thus $z_n$ is forced. By induction, the whole sequence $(z_n)$ is unique.
\end{proof}

Let
\[X=(X_0,X_1,X_2,\ldots), \qquad Y=(Y_0,Y_1,Y_2,\ldots)\]
be sequences of indeterminates. For each $n\geq 0$, define the $n$-th Witt polynomial for the prime $p$ by
\[w_n(X)=X_0^{p^n}+pX_1^{p^{n-1}}+\cdots+p^nX_n.\]

We now look for universal polynomials that can describe addition and multiplication in these coordinates.

\begin{theorem}\label{thm:Witt_poly}
For every $n\geq 0$, there exist unique polynomials
\[S_n,P_n\in \mathbb{Z}[X_0,\dots,X_n,Y_0,\dots,Y_n]\]
such that
\[w_n(S_0,\dots,S_n)=w_n(X_0,\dots,X_n)+w_n(Y_0,\dots,Y_n)\]
and
\[w_n(P_0,\dots,P_n)=w_n(X_0,\dots,X_n)\,w_n(Y_0,\dots,Y_n).\]
\end{theorem}

\begin{proof}
We first prove uniqueness.

Recall that
\[w_n(Z_0,\dots,Z_n)=Z_0^{p^n}+pZ_1^{p^{n-1}}+\cdots+p^nZ_n.\]
Thus the variable $Z_n$ appears only in the last term $p^nZ_n$.
Hence, once $S_0,\dots,S_{n-1}$ are known, the identity
\[w_n(S_0,\dots,S_n)=w_n(X)+w_n(Y)\]
determines $S_n$ uniquely.
The same argument applies to $P_n$.
Therefore, if such polynomials exist, they are unique.

We now prove existence.

Let
\[A=\mathbb{Z}[X_0,X_1,X_2,\dots,Y_0,Y_1,Y_2,\dots]\]
and define a ring endomorphism $\varphi:A\to A$ by
\[\varphi(X_i)=X_i^p,\qquad \varphi(Y_i)=Y_i^p\qquad \text{for all } i\ge 0.\]
Since $\varphi$ is a ring homomorphism and sends each variable to its $p$-th power, one has
\[\varphi(f)\equiv f^p \pmod{pA}\qquad \text{for every } f\in A.\]

Define two sequences $(u_n)_{n\ge0}$ and $(v_n)_{n\ge0}$ in $A$ by
\[u_n:=w_n(X)+w_n(Y),\qquad v_n:=w_n(X)w_n(Y).\]

We claim that these sequences satisfy the congruences of Lemma \ref{lemma:frobenius_congruence}.

First, for $(u_n)$:
\[\varphi(u_n)=\varphi(w_n(X))+\varphi(w_n(Y)).\]
Now
\[\varphi(w_n(X))=X_0^{p^{n+1}}+pX_1^{p^n}+\cdots+p^nX_n^p\equiv w_{n+1}(X)\pmod{p^{n+1}A},\]
since
\[w_{n+1}(X)=X_0^{p^{n+1}}+pX_1^{p^n}+\cdots+p^nX_n^p+p^{n+1}X_{n+1}.\]
Similarly,
\[\varphi(w_n(Y))\equiv w_{n+1}(Y)\pmod{p^{n+1}A}.\]
Therefore
\[\varphi(u_n)\equiv u_{n+1}\pmod{p^{n+1}A}.\]

Next, for $(v_n)$:
\[\varphi(v_n)=\varphi(w_n(X))\,\varphi(w_n(Y)).\]
Using again
\[\varphi(w_n(X))\equiv w_{n+1}(X)\pmod{p^{n+1}A},\qquad\varphi(w_n(Y))\equiv w_{n+1}(Y)\pmod{p^{n+1}A},\]
we obtain
\[\varphi(v_n)\equiv w_{n+1}(X)w_{n+1}(Y)=v_{n+1}\pmod{p^{n+1}A}.\]

Thus both sequences satisfy the hypotheses of Lemma \ref{lemma:frobenius_congruence}. Applying the lemma to $(u_n)$, there exists a unique sequence
\[(S_n)_{n\ge0}\subset A\]
such that
\[u_n=S_0^{p^n}+pS_1^{p^{n-1}}+\cdots+p^nS_n=w_n(S_0,\dots,S_n).\]
That is,
\[w_n(S_0,\dots,S_n)=w_n(X)+w_n(Y).\]

Applying the lemma to $(v_n)$, there exists a unique sequence
\[(P_n)_{n\ge0}\subset A\]
such that
\[v_n=P_0^{p^n}+pP_1^{p^{n-1}}+\cdots+p^nP_n=w_n(P_0,\dots,P_n).\]
That is,
\[w_n(P_0,\dots,P_n)=w_n(X)w_n(Y).\]

Finally, since $u_n$ and $v_n$ involve only the variables
\[X_0,\dots,X_n,Y_0,\dots,Y_n,\]
the recursive construction shows that $S_n$ and $P_n$ also involve only these variables. Hence
\[S_n,P_n\in \mathbb{Z}[X_0,\dots,X_n,Y_0,\dots,Y_n].\]
This completes the proof.
\end{proof}


















\subsection{Related Theorems}


For reasons that will be developed in the next section, we recall several
results explaining why the Witt polynomials can be computed once over
$\mathbb{Z}$ and then reused over any coefficient ring. First we establish
the existence of the canonical morphism $\mathbb{Z}\to A$, then recall the
universal property of polynomial rings, and finally prove the functoriality
of the Witt vector construction.

\begin{proposition}\label{prop:Z_to_A}
Let $A$ be a unital ring. Then there exists a unique unital ring homomorphism
\[\iota_A:\mathbb{Z}\to A.\]
It is given by \[\iota_A(n)=n\cdot 1_A. \]
\end{proposition}

\begin{proof}
Define \[\iota_A(n)=n\cdot 1_A\]
for all $n\in\mathbb{Z}$.

We first check that $\iota_A$ is a ring homomorphism. For $m,n\in\mathbb{Z}$,
\[\iota_A(m+n)=(m+n)1_A=m1_A+n1_A=\iota_A(m)+\iota_A(n).\]

Similarly, \[\iota_A(mn)=(mn)1_A=(m1_A)(n1_A)=\iota_A(m)\iota_A(n).\]

Finally, \[\iota_A(1)=1_A,\]
so $\iota_A$ is unital.

To prove uniqueness, let $f:\mathbb{Z}\to A$ be any unital ring homomorphism.
Since $f(1)=1_A$, for every integer $n\ge0$ we obtain
\[f(n)=f(\underbrace{1+\cdots+1}_{n\text{ times}})=\underbrace{f(1)+\cdots+f(1)}_{n\text{ times}} =n1_A=\iota_A(n).\]

For negative integers, \[f(-n)=-f(n)=-n1_A=\iota_A(-n).\]

Thus $f=\iota_A$, proving uniqueness.
\end{proof}

\begin{theorem}[Extension of coefficient morphisms]
\label{thm:poly_extension}
Let $A$ and $B$ be commutative rings and let \[\varphi:A\to B\]
be a ring homomorphism.

Then there exists a unique ring homomorphism \[\Phi:A[X_0,\dots,X_m]\to B[X_0,\dots,X_m] \]
such that \[\Phi(a)=\varphi(a)\qquad\text{for all }a\in A \]
and \[\Phi(X_i)=X_i\qquad\text{for }i=0,\dots,m.\]

In other words, $\Phi$ applies $\varphi$ to the coefficients of a polynomial
while leaving the variables unchanged.
\end{theorem}

\begin{proof}

Every polynomial in $A[X_0,\dots,X_m]$ can be written uniquely in the form
\[f=\sum_{\alpha} a_\alpha X^\alpha,\]

where the sum runs over finitely many multi-indices \[\alpha=(\alpha_0,\dots,\alpha_m)\in\mathbb{N}^{m+1}.\]

The notation$ X^\alpha$ denotes the monomial \[X^\alpha=X_0^{\alpha_0}X_1^{\alpha_1}\cdots X_m^{\alpha_m}.\]

For example, if \[\alpha=(2,0,3), \]
then \[X^\alpha = X_0^2 X_2^3.\]

Define a map \[\Phi:A[X_0,\dots,X_m]\to B[X_0,\dots,X_m]\]
by \[\Phi\!\left(\sum_{\alpha} a_\alpha X^\alpha\right)=\sum_{\alpha} \varphi(a_\alpha) X^\alpha. \]

Since each polynomial contains only finitely many monomials, the right-hand
side is a well-defined element of $B[X_0,\dots,X_m]$.

We now verify that $\Phi$ is a ring homomorphism.

\textbf{Additivity.}

Let \[f=\sum a_\alpha X^\alpha,\qquad g=\sum b_\alpha X^\alpha.\]

Then \[f+g=\sum (a_\alpha+b_\alpha)X^\alpha.\]

Applying $\Phi$ gives \[\Phi(f+g)=\sum \varphi(a_\alpha+b_\alpha)X^\alpha=\sum (\varphi(a_\alpha)+\varphi(b_\alpha))X^\alpha=\Phi(f)+\Phi(g),\]

since $\varphi$ is a ring homomorphism.

\textbf{Multiplicativity.}

Using the usual multiplication of polynomials, \[fg=\sum_{\alpha,\beta} a_\alpha b_\beta X^{\alpha+\beta},\]

where \[\alpha+\beta=(\alpha_0+\beta_0,\dots,\alpha_m+\beta_m).\]

Applying $\Phi$ gives

\[\Phi(fg)=\sum_{\alpha,\beta}\varphi(a_\alpha b_\beta) X^{\alpha+\beta}.\]

Since $\varphi$ is multiplicative, \[\varphi(a_\alpha b_\beta)=\varphi(a_\alpha)\varphi(b_\beta),\]

so \[\Phi(fg)=\sum_{\alpha,\beta}\varphi(a_\alpha)\varphi(b_\beta) X^{\alpha+\beta}. \]

But this is exactly the product \[\Phi(f)\Phi(g).\]
Thus $\Phi$ preserves multiplication.

Finally, by construction \[\Phi(a)=\varphi(a)\qquad (a\in A), \]

and \[\Phi(X_i)=X_i.\]

To prove uniqueness, observe that every polynomial is a finite sum of
monomials $a_\alpha X^\alpha$. Any ring homomorphism satisfying the above
conditions must therefore send such a monomial to \[\varphi(a_\alpha)X^\alpha. \]

Hence it must coincide with $\Phi$ on all polynomials, and therefore $\Phi$
is unique.

\end{proof}

\begin{theorem}[Functoriality of Witt vectors]
\label{thm:Witt_functoriality}
Let $A$ and $B$ be commutative rings and let \[\varphi:A\to B \]
be a ring homomorphism. Then there exists a unique ring homomorphism \[W(\varphi):W(A)\to W(B) \]
defined by \[W(\varphi)(a_0,a_1,a_2,\dots) =(\varphi(a_0),\varphi(a_1),\varphi(a_2),\dots). \]
\end{theorem}

\begin{proof}
Let \[a=(a_0,a_1,a_2,\dots),\qquad b=(b_0,b_1,b_2,\dots) \]
be Witt vectors in $W(A)$.

By construction of the Witt vector ring, the addition and multiplication
coordinates are given by universal polynomials \[S_n,P_n\in\mathbb{Z}[X_0,\dots,X_n,Y_0,\dots,Y_n] \]
such that \[(a+b)_n=S_n(a_0,\dots,a_n,b_0,\dots,b_n), \]
and \[(ab)_n=P_n(a_0,\dots,a_n,b_0,\dots,b_n).\]

Applying $\varphi$ to the $n$th coordinate of the sum gives \[\varphi((a+b)_n)=\varphi(S_n(a_0,\dots,a_n,b_0,\dots,b_n)).\]

Since $S_n$ has coefficients in $\mathbb{Z}$ and $\varphi$ is a ring
homomorphism, evaluation commutes with $\varphi$, so \[\varphi(S_n(a_0,\dots,a_n,b_0,\dots,b_n))=S_n(\varphi(a_0),\dots,\varphi(a_n),\varphi(b_0),\dots,\varphi(b_n)).\]

Thus \[W(\varphi)(a+b)_n=(W(\varphi)(a)+W(\varphi)(b))_n.\]

The same argument applies to multiplication:
\[\varphi((ab)_n)=\varphi(P_n(a_0,\dots,a_n,b_0,\dots,b_n))=P_n(\varphi(a_0),\dots,\varphi(a_n),\varphi(b_0),\dots,\varphi(b_n)). \]

Hence \[W(\varphi)(ab)=W(\varphi)(a)\,W(\varphi)(b).\]

Therefore $W(\varphi)$ is a ring homomorphism.

Uniqueness follows immediately, since the formula for $W(\varphi)$ is forced
on each coordinate.
\end{proof}

\begin{corollary}\label{cor:ZZ_pipeline}
Let $A$ be a commutative ring. Then the canonical morphism \[\iota_A:\mathbb{Z}\to A\]
induces, by Theorem \ref{thm:poly_extension}, a coefficient morphism \[\mathbb{Z}[X_1,\dots,X_m]\to A[X_1,\dots,X_m].\]

In particular, any polynomial with coefficients in $\mathbb{Z}$ can be
transported canonically to a polynomial with coefficients in $A$.

Moreover, by Theorem \ref{thm:Witt_functoriality}, the morphism $\iota_A$
induces a ring homomorphism \[W(\mathbb{Z})\to W(A). \]

Therefore the universal Witt addition, multiplication, and Frobenius
polynomials, once computed over $\mathbb{Z}$, may be evaluated over any
coefficient ring $A$.
\end{corollary}

\begin{proof}
The first statement follows directly from the universal property of
polynomial rings applied to the canonical morphism \[\iota_A:\mathbb{Z}\to A.\]

The second statement follows from the functoriality of Witt vectors. Since
the Witt sum, product, and Frobenius polynomials have coefficients in
$\mathbb{Z}$, the coefficient morphism \[ \mathbb{Z}[X_1,\dots,X_m]\to A[X_1,\dots,X_m] \]
allows them to be interpreted over $A$, and the induced morphism \[W(\mathbb{Z})\to W(A)\]
shows that this transport is compatible with the Witt vector structure.
\end{proof}

\section{SageMath}
\subsection{Introduction to the software}
SageMath\cite{sagemath}, also called Sage, is an open-source software designed for advanced mathematical computation. It was created over twenty years ago with the goal of providing a free mathematical software environment that integrates many different areas of mathematics, including algebra, geometry, number theory, and topology. By combining numerous specialized mathematical libraries into a single framework, SageMath allows users to access powerful computational tools within a unified interface. Its open-source nature also enables researchers and contributors from around the world to inspect, improve, and extend the software.


\subsection{Current implementation of Witt vectors}

The current implementation of Witt vectors in SageMath \cite{sagemath} uses several different algorithms depending on structural properties of the coefficient ring. These choices reflect mathematical simplifications that occur under certain assumptions on the base ring.

The following pseudo-code summarizes the algorithm selection logic used when constructing a Witt vector ring.

\begin{algorithm}[H]
\caption{Selection of the computation algorithm in a Witt vector ring}
\label{alg:witt_algorithm_selection}
\begin{algorithmic}[1]
\Require Prime $p$, base ring $R$
\If{$p = \operatorname{char}(R)$}
    \If{$R$ is a finite field, or $R$ is a polynomial ring over a finite field}
        \State algorithm $\gets$ \texttt{phantom}
    \Else
        \State algorithm $\gets$ \texttt{finotti}
    \EndIf
\ElsIf{$p$ is a unit in $R$}
    \State algorithm $\gets$ \texttt{p\_invertible}
\Else
    \State algorithm $\gets$ \texttt{standard}
\EndIf
\end{algorithmic}
\end{algorithm}

Each algorithm exploits different algebraic properties of the coefficient ring.

The \emph{phantom} algorithm computes arithmetic in terms of phantom components, which correspond to ghost coordinates lifted to a more convenient ring. Computations are performed in this auxiliary representation and the Witt vector coordinates are reconstructed only when needed.

The \emph{finotti} algorithm applies when the coefficient ring has characteristic $p$. In this case the Witt vector arithmetic can be simplified significantly, and the implementation precomputes binomial tables that accelerate repeated computations.

The \emph{$p$\_invertible} algorithm is used when the prime $p$ is invertible in the coefficient ring. In this situation the ghost component map becomes invertible, allowing Witt vector coordinates to be recovered directly from ghost components without relying on universal Witt polynomials.

Finally, the \emph{standard} algorithm applies in the general case when none of the above simplifications are available. In this setting the Witt vector ring operations are defined through the universal Witt addition and multiplication polynomials introduced in the previous section. These polynomials grow rapidly in complexity as the precision increases, making efficient reuse of previously computed polynomials crucial for performance.

The present work focuses primarily on this \emph{standard} algorithm, since it is the case where the computation and storage of universal Witt polynomials become the main performance bottleneck.


\subsection{UniqueRepresentation}

The current implementation of the class \texttt{WittVectorRing} relies on a Sage-specific mechanism called \emph{UniqueRepresentation}. This mechanism ensures that mathematical parent objects are constructed in a canonical way: if the same object is created multiple times using identical parameters, Sage returns the same instance in memory instead of constructing independent objects.

Internally, \emph{UniqueRepresentation} combines two behaviors provided by the classes \emph{CachedRepresentation} and \emph{WithEqualityById}.

The first mechanism, \emph{CachedRepresentation}, implements a cache of previously constructed objects. When a class inherits from it, Sage records each constructed instance together with the parameters used to create it. If the constructor is called again with the same parameters, the cached instance is returned instead of creating a new object.

For example:

\begin{verbatim}
sage: A = SomeClass(5)
sage: B = SomeClass(5)
sage: A is B
True
\end{verbatim}

The second mechanism, \emph{WithEqualityById}, defines equality between objects purely in terms of identity. Two objects are considered equal if and only if they are the same object in memory.

This approach has important advantages for mathematical parent objects. In particular, it ensures consistent coercion behavior between algebraic structures and avoids the creation of redundant copies of identical rings.

In the case of Witt vector rings, the parameters determining the parent object are the coefficient ring, the precision of the truncated Witt vectors, the prime $p$, and the algorithm used to compute the ring operations. When a Witt vector ring is constructed, Sage checks whether an object with these parameters already exists in the cache. If so, the cached object is returned; otherwise, a new instance is created.


\subsection{Limitations of the current caching approach}

While \emph{UniqueRepresentation} provides canonical construction of parent objects, it ties the caching mechanism directly to the identity of those objects. This design works well when the cached data belongs uniquely to a single mathematical structure.

However, in the case of Witt vectors, a significant portion of the computational cost arises from auxiliary data that is shared across many different Witt vector rings.

For example:

\begin{itemize}
\item Universal Witt addition and multiplication polynomials depend only on the prime $p$.
\item Binomial tables used by the \texttt{finotti} algorithm also depend only on $p$.
\item These objects may be required by many Witt vector rings with different coefficient rings or different precisions.
\end{itemize}

Consequently, two Witt vector rings such as $W_n(R)$ and $W_m(S)$ may be mathematically distinct parent objects, yet still require access to the same intermediate computational data.

Under the \emph{UniqueRepresentation} approach, these auxiliary structures are attached to individual parent objects. As a result, they may be recomputed multiple times across different Witt vector rings even though the mathematical data is identical.

This mismatch between \emph{object identity} and \emph{shared computational data} motivates the use of a different construction mechanism.


\subsection{UniqueFactory}

To address this limitation, Sage provides an alternative construction mechanism called \emph{UniqueFactory}. Instead of attaching the caching mechanism directly to the class representing the mathematical object, a \emph{UniqueFactory} introduces a separate factory responsible for constructing and caching instances.

A typical \emph{UniqueFactory} defines two key methods:

\begin{itemize}
\item \texttt{create\_key}, which validates and normalizes the constructor arguments and produces a canonical key identifying the object,
\item \texttt{create\_object}, which constructs the corresponding instance if it is not already present in the cache.
\end{itemize}

This design separates three responsibilities that are intertwined in \emph{UniqueRepresentation}:

\begin{itemize}
\item normalization of constructor arguments,
\item caching of constructed objects,
\item creation of the underlying mathematical object.
\end{itemize}

Because the factory object exists independently of the mathematical parent class, it can also maintain additional caches of shared computational data. These caches may be indexed differently from the object cache itself.

In the context of Witt vectors, this allows universal Witt polynomials and binomial tables to be stored in caches indexed only by the prime $p$. As a result, these expensive intermediate structures are computed once and reused across many Witt vector rings, even when their coefficient rings or precisions differ.


\subsection{Implementation of the optimizations}
Based on the observations above, the Witt vector implementation was redesigned using a \emph{UniqueFactory}. The implementation of these optimizations is done\cite{grossard_amin_sage_witt_2026}, and a pull request to the official SageMath repository \cite{sagemath} is on the way. In the new design, the public constructor \texttt{WittVectorRing} is no longer a class inheriting from \emph{UniqueRepresentation}, but a factory responsible for constructing and caching the appropriate parent objects.

The following pseudo-code summarizes the construction process.

\begin{algorithm}[H]
\caption{Construction of a truncated Witt vector ring through the factory}
\begin{algorithmic}[1]
\Require coefficient ring $R$, precision $n$, optional prime $p$, optional algorithm $a$
\Ensure A unique Witt vector ring object $W_n(R)$

\State Validate that $R$ is a commutative ring
\State Validate that $n$ is a positive integer
\State Compute $\mathrm{char}(R)$

\If{$p$ is not given}
    \If{$\mathrm{char}(R)$ is prime}
        \State $p \gets \mathrm{char}(R)$
    \Else
        \State raise error
    \EndIf
\ElsIf{$p$ is not prime}
    \State raise error
\EndIf

\State Determine algorithm $a$ using Algorithm~\ref{alg:witt_algorithm_selection}

\State Define normalized key $K = (R,n,p,a)$

\If{$K$ exists in factory cache}
    \State return cached object
\Else
    \State construct algorithm-specific Witt vector ring
    \State store object in factory cache
    \State return object
\EndIf
\end{algorithmic}
\end{algorithm}


\subsection{Shared computational caches}

In addition to the object-level cache, the factory maintains shared caches of auxiliary computational data indexed only by the prime $p$.

For example:

\begin{itemize}
\item binomial tables used by the \texttt{finotti} algorithm,
\item universal Witt polynomials used by the \texttt{standard} algorithm.
\end{itemize}

Before constructing a new Witt vector ring object, the factory therefore ensures that the corresponding shared data is available.

\begin{algorithm}[H]
\caption{Factory dispatch and precomputation}
\begin{algorithmic}[1]
\Require normalized key $K=(R,n,p,a)$

\If{$a=\texttt{finotti}$}
    \If{binomial table for $p$ does not exist}
        \State compute binomial table
    \ElsIf{table too short}
        \State extend table
    \EndIf
    \State create \texttt{WittVectorRing\_finotti}

\ElsIf{$a=\texttt{phantom}$}
    \State create \texttt{WittVectorRing\_phantom}

\ElsIf{$a=\texttt{p\_invertible}$}
    \State create \texttt{WittVectorRing\_pinvertible}

\Else
    \If{universal Witt polynomials for $p$ do not exist}
        \State compute them
    \ElsIf{existing list too short}
        \State extend list
    \EndIf
    \State create \texttt{WittVectorRing\_standard}
\EndIf
\end{algorithmic}
\end{algorithm}


\subsection{Use of universal Witt polynomials}

For the \texttt{standard} algorithm, the key observation is that the Witt addition, multiplication, and Frobenius polynomials are universal polynomials defined over $\mathbb{Z}$. As shown in the previous section, these polynomials depend only on the prime $p$ and not on the coefficient ring.

Consequently, the factory computes these polynomials once over $\mathbb{Z}$ and stores them in its cache indexed by $p$. When a specific Witt vector ring over a coefficient ring $R$ is created, the factory converts these universal polynomials to the corresponding polynomial ring over $R$.

\begin{algorithm}[H]
\caption{Use of cached universal polynomials}
\begin{algorithmic}[1]
\Require coefficient ring $R$, precision $n$, prime $p$

\State retrieve universal Witt polynomials for $p$

\For{$i=0,\dots,n-1$}
    \State convert $i$-th universal sum polynomial to ring over $R$
    \State convert $i$-th universal product polynomial to ring over $R$
\EndFor

\For{$i=0,\dots,n-2$}
    \State convert $i$-th Frobenius polynomial to ring over $R$
\EndFor

\State attach converted polynomials to the Witt vector ring object
\end{algorithmic}
\end{algorithm}

\section{Benchmarks}
Yay it helps a lot

\section{Other things i learned}
\subsection{Git}
Git branching and rebasing/merging, PR
\subsection{Software engineering practices}
Python PEP8, Docstring/Doctests and coverage, readability > micro optimizations.

\newpage


\newpage
\bibliographystyle{plain}
\bibliography{main}

\newpage
\appendix

\section{Appendix}
\label{sec:appendix}
Proof of \ref{equiv_A}.
\begin{proof}
  We will proceed by recursion on $i\in\mathbb{N}$, with the case $i=0$ being the hypothesis.
  Suppose now that for $i\in\mathbb{N}$ we have $x^{p^i}\equiv y^{p^i} \pmod {p^{i+1}A}$. We can write, for some $\alpha \in A $ \[x^{p^i}=y^{p^i}+\alpha p^{i+1}\implies x^{p^{i+1}}=\bigl(y^{p^i}+\alpha p^{i+1}\bigr)^p.\]
  Now expanding with Newton's binomial theorem, \[x^{p^{i+1}}=\bigl(y^{p^(i+1)}+\alpha p^{i+1}\bigr)^p=y^{p^i}+\alpha^p p^{p(i+1)}+\sum_{k=0}^{p-1}\binom{p}{k}y^{p^{i}(p-k)}\alpha^k p^{k(i+1)}.\]
  For $k\in\{1,\ldots,p-1\}$ we have \[\binom{p}{k}\mid pA \implies \binom{p}{k}y^{p^{i}(p-k)}\alpha^k p^{k(i+1)}\mid p^{i+2}.\]
  And \[p^{p(i+1)}\mid p^{i+2} \text{ since }p\geq 2.\]
  We indeed get that $x^{p^{i+1}}\equiv y^{p^{i+1}} \pmod {p^{i+2}A}$.
\end{proof}
\end{document}
